\documentclass{lecture-kuznia}

\def\data{26.09.2026}
\def\place{Poręba Wielka}
\def\title{Od AM--GM do AM--HM i nie tylko}
\def\group{Juniorzy}
\def\author{Autor: Stefan Świerczewski}
\def\lecturer{Prowadzący: Stefan Świerczewski}

\begin{document}

\begin{center}
	{\LARGE\textbf{\title{}}}
\end{center}

\section{Teoria}

\subsection{Cztery średnie}

Dla dodatnich liczb $a_1,\dots,a_n$ definiujemy
\[
	A=\frac{a_1+\cdots+a_n}{n},\qquad
	G=\sqrt[n]{a_1\cdots a_n},
\]
\[
	H=\frac{n}{\frac1{a_1}+\cdots+\frac1{a_n}},\qquad
	Q=\sqrt{\frac{a_1^2+\cdots+a_n^2}{n}}.
\]

Średnią $Q$ nazywamy \emph{średnią kwadratową}; spotyka się też oznaczenie
$P_2$.

Istnieje również szersza rodzina tzw.\ \emph{średnich potęgowych},
obejmująca między innymi średnią harmoniczną, arytmetyczną i kwadratową.
Nie będziemy tutaj rozwijać ich ogólnej teorii.

Najważniejszy łańcuch ma postać
\[
	\boxed{Q\ge A\ge G\ge H}.
\]

Równość w którejkolwiek z tych nierówności zachodzi wtedy i tylko wtedy,
gdy
\[
	a_1=a_2=\cdots=a_n.
\]

\subsection{Dlaczego \texorpdfstring{$Q\ge A$}{Q ≥ A}?}

Mamy
\[
	n(a_1^2+\cdots+a_n^2)-(a_1+\cdots+a_n)^2
	=
	\sum_{1\le i<j\le n}(a_i-a_j)^2\ge0.
\]

Zatem
\[
	(a_1+\cdots+a_n)^2\le n(a_1^2+\cdots+a_n^2),
\]
czyli
\[
	\frac{a_1+\cdots+a_n}{n}
	\le
	\sqrt{\frac{a_1^2+\cdots+a_n^2}{n}}.
\]

\subsection{AM--GM}

Dla $a,b>0$ mamy
\[
	\frac{a+b}{2}\ge \sqrt{ab},
\]
bo
\[
	(\sqrt a-\sqrt b)^2\ge0.
\]

W ogólności, dla $a_1,\dots,a_n>0$,
\[
	\frac{a_1+\cdots+a_n}{n}\ge \sqrt[n]{a_1\cdots a_n}.
\]

\begin{proof}
	Najpierw udowodnimy pomocniczy fakt: jeśli
	\[
		a_1a_2\cdots a_n=1,
	\]
	to
	\[
		a_1+\cdots+a_n\ge n,
	\]
	przy czym równość zachodzi wtedy i tylko wtedy, gdy
	\[
		a_1=\cdots=a_n=1.
	\]

	Dowód przeprowadzimy indukcyjnie względem $n$.
	Dla $n=1$ teza jest oczywista.

	Załóżmy więc, że $n\ge2$ i że teza zachodzi dla $n-1$ liczb.
	Bez straty ogólności możemy przyjąć, że
	\[
		a_1=\max\{a_1,\dots,a_n\}\ge1,
		\qquad
		a_2=\min\{a_1,\dots,a_n\}\le1.
	\]

	Wtedy
	\[
		a_1+a_2-a_1a_2-1
		=
		(a_1-1)(1-a_2)\ge0,
	\]
	czyli
	\[
		a_1+a_2-a_1a_2\ge1.
	\]

	Ponadto liczby
	\[
		a_1a_2,\ a_3,\dots,a_n
	\]
	mają iloczyn równy $1$, więc z założenia indukcyjnego
	\[
		a_1a_2+a_3+\cdots+a_n\ge n-1.
	\]

	Dodając obie nierówności, otrzymujemy
	\[
		a_1+a_2+\cdots+a_n\ge n.
	\]

	Jeśli zachodzi równość, to w szczególności
	\[
		(a_1-1)(1-a_2)=0.
	\]

	Stąd $a_1=1$ lub $a_2=1$. Ponieważ $a_1$ jest największą, a $a_2$
	najmniejszą spośród liczb $a_i$, w obu przypadkach otrzymujemy
	\[
		a_1=\cdots=a_n=1.
	\]

	Teraz udowodnimy AM--GM. Niech
	\[
		g=\sqrt[n]{a_1a_2\cdots a_n}.
	\]

	Liczby
	\[
		\frac{a_1}{g},\frac{a_2}{g},\dots,\frac{a_n}{g}
	\]
	mają iloczyn równy $1$. Z udowodnionego wyżej faktu
	\[
		\frac{a_1}{g}+\frac{a_2}{g}+\cdots+\frac{a_n}{g}\ge n.
	\]

	Po pomnożeniu przez $g$ dostajemy
	\[
		a_1+a_2+\cdots+a_n\ge ng,
	\]
	czyli
	\[
		\frac{a_1+\cdots+a_n}{n}
		\ge
		\sqrt[n]{a_1a_2\cdots a_n}.
	\]

	Równość zachodzi wtedy i tylko wtedy, gdy
	\[
		\frac{a_1}{g}=\frac{a_2}{g}=\cdots=\frac{a_n}{g}=1,
	\]
	czyli
	\[
		a_1=a_2=\cdots=a_n.
	\]
\end{proof}

\begin{example}\label{E1}
	Wyznacz minimum
	\[
		x+\frac{4}{x^2},\qquad x>0.
	\]
\end{example}

\begin{solution*}[\ref{E1}]
	Rozbijamy pierwszy składnik:
	\[
		x+\frac4{x^2}
		=
		\frac x2+\frac x2+\frac4{x^2}.
	\]

	Iloczyn tych trzech liczb jest równy $1$, więc
	\[
		x+\frac4{x^2}\ge3.
	\]

	Równość zachodzi dla
	\[
		\frac x2=\frac4{x^2},
	\]
	czyli dla $x=2$.
\end{solution*}

\subsection{Jak dobierać składniki?}

Sztuczka często polega nie na samym użyciu AM--GM,
lecz na takim zapisaniu wyrażenia, aby:

\begin{itemize}
	\item wszystkie składniki były dodatnie,
	\item ich suma albo iloczyn był łatwy do kontrolowania
	(o kontrolowaniu uchylimy rąbka tajemnicy za chwilę),
	\item warunek równości dało się rzeczywiście spełnić.
\end{itemize}

Na przykład, aby oszacować $x(4-x)^2$, naturalnie patrzymy na
\[
	x,\qquad \frac{4-x}{2},\qquad \frac{4-x}{2},
\]
bo ich suma jest równa $4$.

Jeżeli potęgi są różne, powtarzamy składniki. Wyrażenie
\[
	xy^2z^3
\]
sugeruje sześć liczb
\[
	x,\quad \frac y2,\frac y2,\quad
	\frac z3,\frac z3,\frac z3.
\]

\subsection{Jednorodność i normalizacja}

Nierówność nazywamy \emph{jednorodną stopnia $d$}, jeżeli po zastąpieniu
\[
	(a,b,c)\longmapsto(ta,tb,tc)
\]
obie strony mnożą się przez $t^d$.

Przykładowo
\[
	(a+b+c)^2\ge3(ab+bc+ca)
\]
jest jednorodna stopnia $2$, a
\[
	(a+b+c)\left(\frac1a+\frac1b+\frac1c\right)\ge9
\]
jest jednorodna stopnia $0$: po przeskalowaniu zmiennych całe wyrażenie
w ogóle się nie zmienia.

Jednorodność pozwala nam \emph{znormalizować} zmienne. Jeśli nierówność
jest jednorodna, możemy często bez straty ogólności założyć na przykład
\[
	a+b+c=1,\qquad a+b+c=3
\]
albo
\[
	abc=1,
\]
wybierając odpowiedni wspólny mnożnik.

Czasem robimy operację odwrotną: \emph{ujednoradniamy} zadanie.
Na przykład przy warunku $a+b+c=3$ nierówność
\[
	a^2+b^2+c^2\ge3
\]
wygodniej zapisać jako
\[
	3(a^2+b^2+c^2)\ge(a+b+c)^2.
\]

Stała zniknęła, a otrzymana nierówność jest jednorodna.

\begin{remark}
	Normalizować można tylko wtedy, gdy skala rzeczywiście nie ma znaczenia.
	Jeżeli po pomnożeniu wszystkich zmiennych przez $t$ obie strony zmieniają
	się inaczej, nie wolno po prostu założyć $a+b+c=1$.
\end{remark}

\subsection{AM--HM}

Stosując AM--GM do liczb odwrotnych, otrzymujemy
\[
	\frac{\frac1{a_1}+\cdots+\frac1{a_n}}{n}
	\ge
	\frac1{\sqrt[n]{a_1\cdots a_n}},
\]
a po odwróceniu
\[
	H\le G.
\]

AM--HM jest szczególnie wygodne wtedy, gdy w zadaniu pojawia się suma
odwrotności.

\begin{example}\label{E2}
	Dla $a,b,c>0$ pokaż, że
	\[
		\frac1{a+b}+\frac1{b+c}+\frac1{c+a}
		\ge
		\frac9{2(a+b+c)}.
	\]
\end{example}

\begin{solution*}[\ref{E2}]
	Stosujemy AM--HM do liczb
	\[
		a+b,\qquad b+c,\qquad c+a.
	\]

	Ich średnia arytmetyczna wynosi
	\[
		\frac{2(a+b+c)}3,
	\]
	więc
	\[
		\frac{3}{
			\frac1{a+b}+\frac1{b+c}+\frac1{c+a}
		}
		\le
		\frac{2(a+b+c)}3,
	\]
	co po przekształceniu daje tezę.
\end{solution*}

\newpage

\section{Zadania}

\begin{problem}{$\clubsuit$}\label{Z1}
	Wyznacz najmniejszą wartość
	\[
		(x+3)\left(1+\frac3x\right),
		\qquad x>0.
	\]
	Dla jakiego $x$ jest ona osiągana?
\end{problem}

\begin{problem}\label{Z2}
	Niech $a,b,c>0$ oraz $a+b+c=3$. Udowodnij, że
	\[
		(1+a)(1+b)(1+c)\le8.
	\]
	Kiedy zachodzi równość?
\end{problem}

\begin{problem}\label{Z3}
	Niech $a,b,c>0$ oraz $abc=1$. Udowodnij, że
	\[
		(1+a)(1+b)(1+c)\ge8.
	\]
	Porównaj pomysł z poprzednim zadaniem.
\end{problem}

\begin{problem}\label{Z4}
	Liczby $a,b,c$ są dodatnie i $abc=1$. Udowodnij, że
	\[
		(a+b)(b+c)(c+a)\ge8.
	\]
\end{problem}

\begin{problem}{$\diamondsuit$}\label{Z5}
	Dla każdej dodatniej liczby naturalnej $n$ udowodnij
	\[
		n!\le\left(\frac{n+1}{2}\right)^n.
	\]
	Dla jakich $n$ zachodzi równość?
\end{problem}

\begin{problem}{$\clubsuit$}\label{Z6}
	Prostopadłościan ma dodatnie krawędzie $a,b,c$ i pole powierzchni równe $96$.
	Wyznacz największą możliwą objętość prostopadłościanu.
\end{problem}

\begin{problem}[$\redheart$]\label{Z7}
	Dla dodatnich $a,b,c$ udowodnij
	\[
		(a+b+c)^2\ge3(ab+bc+ca).
	\]
\end{problem}

\begin{problem}\label{Z8}
	Dla dodatnich $a,b,c$ wyznacz najmniejszą wartość
	\[
		\frac{a^2+b^2+c^2}{(a+b+c)^2}.
	\]
\end{problem}

\begin{problem}\label{Z9}
	Dla $a,b,c>0$ udowodnij
	\[
		(a+b+c)(ab+bc+ca)\ge9abc.
	\]
\end{problem}

\begin{problem}{$\diamondsuit$}\label{Z10}
	Udowodnij, że dla dowolnych liczb dodatnich $a,b$ oraz liczby naturalnej $m$
	zachodzi nierówność
	\[
		\left(1+\frac ab\right)^m+
		\left(1+\frac ba\right)^m
		\ge 2^{m+1}.
	\]
\end{problem}

\begin{problem}[$\redheart$]\label{Z11}
	Niech $a,b,c>0$ oraz $abc=1$. Udowodnij, że
	\[
		a^2+b^2+c^2\ge a+b+c.
	\]
\end{problem}

\begin{problem}{$\spadesuit$}\label{Z12}
	Udowodnij, że dla dowolnej liczby naturalnej $n\ge2$ zachodzą nierówności
	\[
		\left(\frac{2n+1}{3}\right)^{\frac{n(n+1)}2}
		>
		\prod_{i=2}^{n} i^i
		>
		\left(\frac{n+1}{2}\right)^{\frac{n(n+1)}2}.
	\]
\end{problem}

\begin{problem}{$\spadesuit$}\label{Z13}
	Niech $a_1,a_2,\dots,a_n$ będą dodatnimi liczbami oraz niech
	$q_1,q_2,\dots,q_n$ będą dodatnimi liczbami wymiernymi spełniającymi(uwaga w ogólności wystarczy, że są to liczby dodatnie)
	\[
		q_1+q_2+\cdots+q_n=1.
	\]
	Udowodnij ważoną nierówność AM--GM
	\[
		q_1a_1+q_2a_2+\cdots+q_na_n
		\ge
		a_1^{q_1}a_2^{q_2}\cdots a_n^{q_n}.
	\]
	Udowodnij również, że równość zachodzi wtedy i tylko wtedy, gdy
	\[
		a_1=a_2=\cdots=a_n.
	\]
\end{problem}

\begin{problem}\label{Z14}
	Niech $a,b,c>0$ oraz
	\[
		a^2+b^2+c^2=1.
	\]
	Udowodnij, że
	\[
		\frac{ab}{c}+\frac{bc}{a}+\frac{ca}{b}\ge\sqrt3.
	\]
	Kiedy zachodzi równość?
\end{problem}

\begin{problem}\label{Z15}
	Niech $a,b,c>0$ oraz
	\[
		a^2+b^2+c^2=1.
	\]
	Udowodnij, że
	\[
		a^3\left(\frac bc+\frac cb\right)
		+b^3\left(\frac ca+\frac ac\right)
		+c^3\left(\frac ab+\frac ba\right)
		+3abc
		\ge\sqrt3.
	\]
\end{problem}

\begin{problem}[$\redheart$]\label{Z16}
	Dla dodatnich $a,b,c$ udowodnij
	\[
		(a^2b+b^2c+c^2a)(ab^2+bc^2+ca^2)
		\ge 9a^2b^2c^2.
	\]
\end{problem}

\begin{problem}[$\redheart$]\label{Z17}
	Wypukły $n$-kąt ma obwód $P$. Udowodnij, że istnieją trzy kolejne
	wierzchołki tworzące trójkąt o polu nie większym niż
	\[
		\frac{P^2}{2n^2}.
	\]
\end{problem}

% \begin{problem}[$\redheart\ \redheart$]\label{Z18}
% Niech $G$ będzie grafem o $n$ wierzchołkach, którego największy zbiór
% niezależny ma rozmiar co najwyżej $\alpha$.

% Krawędzie grafu pokryto pełnymi grafami dwudzielnymi
% \[
% K_{A_1,B_1},K_{A_2,B_2},\dots,K_{A_m,B_m}.
% \]

% To znaczy: każda krawędź $G$ należy do co najmniej jednego z tych grafów.
% Udowodnij, że
% \[
% \sum_{i=1}^m\bigl(|A_i|+|B_i|\bigr)
% \ge
% n\log_2\frac{n}{\alpha}.
% \]
% \end{problem}

\ifshowsolutions

\newpage

\section{Rozwiązania}

\begin{solution*}[\ref{Z1}]
	Po rozwinięciu
	\[
		(x+3)\left(1+\frac3x\right)
		=x+6+\frac9x.
	\]
	Z AM--GM
	\[
		x+\frac9x\ge6,
	\]
	więc całe wyrażenie jest co najmniej równe $12$.
	Równość zachodzi dla $x=3$.
\end{solution*}

\begin{solution*}[\ref{Z2}]
	Liczby
	\[
		1+a,\qquad1+b,\qquad1+c
	\]
	mają sumę równą $6$. Zatem
	\[
		(1+a)(1+b)(1+c)
		\le
		\left(\frac63\right)^3=8.
	\]
	Równość zachodzi dla $a=b=c=1$.
\end{solution*}

\begin{solution*}[\ref{Z3}]
	Dla każdej dodatniej liczby $t$ mamy
	\[
		1+t\ge2\sqrt t.
	\]
	Stąd
	\[
		(1+a)(1+b)(1+c)
		\ge
		8\sqrt{abc}=8.
	\]
	Równość zachodzi dla $a=b=c=1$.
\end{solution*}

\begin{solution*}[\ref{Z4}]
	Mamy
	\[
		a+b\ge2\sqrt{ab},\qquad
		b+c\ge2\sqrt{bc},\qquad
		c+a\ge2\sqrt{ca}.
	\]
	Po pomnożeniu:
	\[
		(a+b)(b+c)(c+a)\ge8abc=8.
	\]
\end{solution*}

\begin{solution*}[\ref{Z5}]
	Średnia arytmetyczna liczb $1,2,\dots,n$ wynosi
	\[
		\frac{1+2+\cdots+n}{n}=\frac{n+1}{2}.
	\]
	Z AM--GM
	\[
		\sqrt[n]{n!}\le\frac{n+1}{2},
	\]
	a więc
	\[
		n!\le\left(\frac{n+1}{2}\right)^n.
	\]
	Równość zachodzi tylko dla $n=1$.
\end{solution*}

\begin{solution*}[\ref{Z6}]
	Z warunku na pole powierzchni
	\[
		2(ab+bc+ca)=96,
	\]
	czyli
	\[
		ab+bc+ca=48.
	\]
	Z AM--GM dla $ab,bc,ca$:
	\[
		16=\frac{ab+bc+ca}{3}
		\ge
		\sqrt[3]{a^2b^2c^2}
		=(abc)^{2/3}.
	\]
	Zatem
	\[
		abc\le64.
	\]
	Równość zachodzi dla $a=b=c=4$.
\end{solution*}

\begin{solution*}[\ref{Z7}]
	Dodając
	\[
		a^2+b^2\ge2ab,\qquad
		b^2+c^2\ge2bc,\qquad
		c^2+a^2\ge2ca,
	\]
	otrzymujemy
	\[
		a^2+b^2+c^2\ge ab+bc+ca.
	\]
	Po dodaniu $2(ab+bc+ca)$ do obu stron:
	\[
		(a+b+c)^2\ge3(ab+bc+ca).
	\]
\end{solution*}

\begin{solution*}[\ref{Z8}]
	Wyrażenie jest jednorodne stopnia $0$, więc możemy założyć
	\[
		a+b+c=1.
	\]
	Z $Q\ge A$:
	\[
		a^2+b^2+c^2
		\ge
		\frac{(a+b+c)^2}{3}
		=
		\frac13.
	\]
	Minimum wynosi $\frac13$ i jest osiągane dla $a=b=c$.
\end{solution*}

\begin{solution*}[\ref{Z9}]
	Z AM--GM
	\[
		a+b+c\ge3\sqrt[3]{abc}
	\]
	oraz
	\[
		ab+bc+ca
		\ge
		3\sqrt[3]{(ab)(bc)(ca)}
		=
		3\sqrt[3]{a^2b^2c^2}.
	\]
	Po pomnożeniu otrzymujemy
	\[
		(a+b+c)(ab+bc+ca)\ge9abc.
	\]
\end{solution*}

\begin{solution*}[\ref{Z10}]
	Oznaczmy
	\[
		x=\frac ab>0.
	\]
	Wtedy
	\[
		\left(1+\frac ab\right)^m+\left(1+\frac ba\right)^m
		=
		(1+x)^m+\left(1+\frac1x\right)^m.
	\]
	Z AM--GM dla dwóch składników:
	\[
		(1+x)^m+\left(1+\frac1x\right)^m
		\ge
		2\sqrt{
			\left((1+x)\left(1+\frac1x\right)\right)^m
		}.
	\]
	Ponieważ
	\[
		(1+x)\left(1+\frac1x\right)
		=
		2+x+\frac1x
		\ge4,
	\]
	otrzymujemy
	\[
		\left(1+\frac ab\right)^m+\left(1+\frac ba\right)^m
		\ge
		2\cdot4^{m/2}
		=
		2^{m+1}.
	\]
	Równość zachodzi dla $a=b$.
\end{solution*}

\begin{solution*}[\ref{Z11}]
	Stosujemy AM--GM do sześciu liczb
	\[
		a^2,a^2,a^2,a^2,b^2,c^2.
	\]
	Dostajemy
	\[
		\frac{4a^2+b^2+c^2}{6}
		\ge
		\sqrt[6]{a^8b^2c^2}.
	\]
	Ponieważ $abc=1$, mamy
	\[
		a^8b^2c^2=a^6,
	\]
	więc
	\[
		4a^2+b^2+c^2\ge6a.
	\]
	Analogicznie
	\[
		a^2+4b^2+c^2\ge6b,
		\qquad
		a^2+b^2+4c^2\ge6c.
	\]
	Po dodaniu:
	\[
		6(a^2+b^2+c^2)\ge6(a+b+c).
	\]
	Zatem
	\[
		a^2+b^2+c^2\ge a+b+c.
	\]
\end{solution*}

\begin{solution*}[\ref{Z12}]
	Rozważmy następujący zbiór
	\[
		\frac{n(n+1)}2
	\]
	liczb:
	\[
		1,\quad
		2,2,\quad
		3,3,3,\quad
		\dots,\quad
		\underbrace{n,n,\dots,n}_{n}.
	\]

	Ich iloczyn wynosi
	\[
		1^1\cdot2^2\cdot3^3\cdots n^n
		=
		\prod_{i=2}^n i^i,
	\]
	natomiast ich suma jest równa
	\[
		1^2+2^2+\cdots+n^2
		=
		\frac{n(n+1)(2n+1)}6.
	\]

	Z AM--GM
	\[
		\sqrt[\frac{n(n+1)}2]
		{\prod_{i=2}^n i^i}
		<
		\frac{
			\frac{n(n+1)(2n+1)}6
		}{
			\frac{n(n+1)}2
		}
		=
		\frac{2n+1}{3}.
	\]

	Stąd
	\[
		\prod_{i=2}^n i^i
		<
		\left(\frac{2n+1}{3}\right)^{\frac{n(n+1)}2}.
	\]

	Dla drugiej nierówności rozważmy liczby
	\[
		1,\quad
		\frac12,\frac12,\quad
		\frac13,\frac13,\frac13,\quad
		\dots,\quad
		\underbrace{\frac1n,\dots,\frac1n}_{n}.
	\]

	Ponownie jest ich $\frac{n(n+1)}2$, a ich suma wynosi
	\[
		1+1+\cdots+1=n.
	\]

	Z AM--GM
	\[
		\sqrt[\frac{n(n+1)}2]
		{
			\frac1{\prod_{i=2}^n i^i}
		}
		<
		\frac{n}{\frac{n(n+1)}2}
		=
		\frac2{n+1}.
	\]

	Zatem
	\[
		\frac1{\prod_{i=2}^n i^i}
		<
		\left(\frac2{n+1}\right)^{\frac{n(n+1)}2},
	\]
	a po odwróceniu
	\[
		\prod_{i=2}^n i^i
		>
		\left(\frac{n+1}{2}\right)^{\frac{n(n+1)}2}.
	\]
\end{solution*}

\begin{solution*}[\ref{Z13}]
	Najpierw załóżmy, że wszystkie wagi $q_i$ są wymierne. Możemy wtedy
	zapisać
	\[
		q_i=\frac{k_i}{m},
	\]
	gdzie $k_1,\dots,k_n$ są dodatnimi liczbami naturalnymi oraz
	\[
		k_1+\cdots+k_n=m.
	\]

	Rozważmy $m$ liczb, wśród których $a_i$ występuje dokładnie $k_i$ razy.
	Z AM--GM otrzymujemy
	\[
		\frac{k_1a_1+\cdots+k_na_n}{m}
		\ge
		\sqrt[m]{a_1^{k_1}\cdots a_n^{k_n}}.
	\]

	Lewa strona jest równa
	\[
		q_1a_1+\cdots+q_na_n,
	\]
	natomiast prawa
	\[
		a_1^{q_1}\cdots a_n^{q_n}.
	\]

	Zatem
	\[
		q_1a_1+\cdots+q_na_n
		\ge
		a_1^{q_1}\cdots a_n^{q_n}.
	\]

	Dla dowolnych rzeczywistych wag dodatnich wynik otrzymujemy przez
	przybliżanie wag liczbami wymiernymi i przejście do granicy.

	Równość w AM--GM zachodzi wtedy i tylko wtedy, gdy wszystkie liczby,
	do których ją stosujemy, są równe. Ponieważ wszystkie wagi są dodatnie,
	oznacza to dokładnie
	\[
		a_1=a_2=\cdots=a_n.
	\]
\end{solution*}

\begin{solution*}[\ref{Z14}]
	Wprowadźmy
	\[
		x=\frac{ab}{c},\qquad
		y=\frac{bc}{a},\qquad
		z=\frac{ca}{b}.
	\]
	Wtedy
	\[
		xy=b^2,\qquad yz=c^2,\qquad zx=a^2,
	\]
	a zatem
	\[
		xy+yz+zx=a^2+b^2+c^2=1.
	\]
	Ponieważ
	\[
		(x+y+z)^2\ge3(xy+yz+zx),
	\]
	mamy
	\[
		x+y+z\ge\sqrt3.
	\]
	Po powrocie do $a,b,c$ dostajemy
	\[
		\frac{ab}{c}+\frac{bc}{a}+\frac{ca}{b}\ge\sqrt3.
	\]
	Równość zachodzi wtedy i tylko wtedy, gdy $x=y=z$, czyli $a=b=c$.
	Z warunku
	\[
		a^2+b^2+c^2=1
	\]
	otrzymujemy wtedy
	\[
		a=b=c=\frac1{\sqrt3}.
	\]
\end{solution*}

\begin{solution*}[\ref{Z15}]
	Lewą stronę można rozpoznać jako iloczyn
	\[
		(a^2+b^2+c^2)
		\left(
		\frac{ab}{c}+\frac{bc}{a}+\frac{ca}{b}
		\right).
	\]
	Rzeczywiście, po wymnożeniu otrzymujemy dokładnie
	\[
		a^3\left(\frac bc+\frac cb\right)
		+b^3\left(\frac ca+\frac ac\right)
		+c^3\left(\frac ab+\frac ba\right)
		+3abc.
	\]
	Ponieważ
	\[
		a^2+b^2+c^2=1,
	\]
	lewa strona jest równa
	\[
		\frac{ab}{c}+\frac{bc}{a}+\frac{ca}{b}.
	\]
	Teza wynika więc bezpośrednio z poprzedniego zadania.
\end{solution*}

\begin{solution*}[\ref{Z16}]
	W każdym nawiasie stosujemy AM--GM:
	\[
		a^2b+b^2c+c^2a
		\ge
		3\sqrt[3]{(a^2b)(b^2c)(c^2a)}
		=
		3abc.
	\]
	Podobnie
	\[
		ab^2+bc^2+ca^2
		\ge3abc.
	\]
	Po pomnożeniu dostajemy
	\[
		(a^2b+b^2c+c^2a)(ab^2+bc^2+ca^2)
		\ge9a^2b^2c^2.
	\]
\end{solution*}

\begin{solution*}[\ref{Z17}]
	Niech długości kolejnych boków będą równe
	\[
		x_1,x_2,\dots,x_n.
	\]
	Suma liczb
	\[
		x_1+x_2,\ x_2+x_3,\ \dots,\ x_n+x_1
	\]
	wynosi $2P$. Z zasady szufladkowej dla pewnego $i$ mamy więc
	\[
		x_i+x_{i+1}\le\frac{2P}{n}.
	\]
	Pole trójkąta wyznaczonego przez te dwa kolejne boki jest nie większe niż
	\[
		\frac12x_ix_{i+1}.
	\]
	Z AM--GM
	\[
		x_ix_{i+1}
		\le
		\left(\frac{x_i+x_{i+1}}2\right)^2.
	\]
	Stąd
	\[
		[\triangle]
		\le
		\frac{(x_i+x_{i+1})^2}{8}
		\le
		\frac18\left(\frac{2P}{n}\right)^2
		=
		\frac{P^2}{2n^2}.
	\]
\end{solution*}

% \begin{solution*}[\ref{Z18}]
% Dla każdego $i$ wybierzmy niezależnie i z prawdopodobieństwem $\frac12$
% jedną z dwóch części $A_i,B_i$.

% Nazwijmy wierzchołek $v$ \emph{dobrym}, jeżeli we wszystkich grafach
% dwudzielnych, do których należy, znajduje się po wybranej stronie.

% Dwa sąsiednie wierzchołki nie mogą być jednocześnie dobre: istnieje bowiem
% jeden z pokrywających grafów dwudzielnych, w którym leżą po przeciwnych
% stronach. Zatem zbiór dobrych wierzchołków jest niezależny, więc zawsze ma
% co najwyżej $\alpha$ elementów.

% Niech $d(v)$ oznacza liczbę grafów dwudzielnych zawierających $v$.
% Wtedy
% \[
% \Pr(v\text{ jest dobry})=2^{-d(v)}.
% \]

% Z liniowości wartości oczekiwanej
% \[
% \mathbb E|X|
% =
% \sum_v2^{-d(v)}.
% \]

% Z AM--GM:
% \[
% \frac1n\sum_v2^{-d(v)}
% \ge
% \sqrt[n]{\prod_v2^{-d(v)}}
% =
% 2^{-\frac1n\sum_vd(v)}.
% \]

% Jeśli
% \[
% S=\sum_{i=1}^m(|A_i|+|B_i|),
% \]
% to
% \[
% S=\sum_vd(v).
% \]

% A zatem
% \[
% \alpha
% \ge
% \mathbb E|X|
% \ge
% n2^{-S/n}.
% \]

% Po przekształceniu
% \[
% S\ge n\log_2\frac{n}{\alpha}.
% \]
% \end{solution*}

\fi

\end{document}
